Ukázalo sa však, že DNA sekvencie v klastroch sa zhodovali iba v úvodnom úseku dlhom niekoľko nukleotidov, čo znamenalo, že klastrovanie nebolo dostatočne presné.
Neskôr sa zistilo, že Clover vykonáva zarovnanie iba podľa začiatku reťazca a neberie do úvahy zhodu v ďalších častiach sekvencie. V jednom klastri sa preto mohli nachádzať sekvencie, ktoré mali zhodu iba na začiatku a ďalej sa výrazne rozchádzali.
Jeden klaster obsahoval viac ako 37~000 sekvencií, čo predstavovalo neobvykle veľký počet čítaní. Spoločná zhoda týchto sekvencií mala dĺžku iba 43 nukleotidov, pričom takýto jav sa javil ako málo pravdepodobný pri náhodnom výskyte.
